Let the smallest angle be $2\theta$, and scale the five sides to one. The extremal convex configuration is approached when the two edges on each side of that angle straighten into length-two legs; the remaining unit edge joins their ends. In that limiting isosceles triangle, $1=4\sin\theta$, so $\sin\theta=\dfrac{1}{4}$. To justify the bound, straighten the two two-edge boundary chains outward while keeping them inside the angle and preserving their unit lengths: convexity makes the distance between their outer ends nondecreasing, with maximum $4\sin\theta$ when both chains are straight. Their original endpoints were one unit apart, hence $1\le4\sin\theta$. The bound is approached by strictly convex perturbations of the straight-chain configuration. Thus $\sin(2\theta)=2(\dfrac{1}{4})(\dfrac{\sqrt{15}}{4})=\boxed{\dfrac{\sqrt{15}}{8}}$.
\begin{asy}
import graph;import geometry;usepackage("amsmath");size(7cm);
real t=asin(.25)*180/pi;pair A=(0,0),B=dir(90+t),C=2*B,E=dir(90-t),D=2*E;draw(A--C--D--A);dot(B);dot(E);label("$A$",A,(0,-1));label("$B$",B,(-1,0));label("$C$",C,(-1,1));label("$D$",D,(1,1));label("$E$",E,(1,0));label("limiting configuration",(0,2.5));
\end{asy}

